\section{Chebyshev improvement of physical viability}

Use the supplied windows, core tensors and closing matrix of
Theorem~\ref{thm:physical_polynomial_tensor_family}. In particular,
$a+2\ell\le b$, $b+2\ell\le L-1$, $\ell\ge1$, $q\ge2$ and $k\ge1$;
$W^L_r,W^R_r$ have the stated two-site supports. The core is the actual trace
contraction (\ref{eq:physical_tensor_core}), with arbitrary closing matrix
and $B\ge0$. Retain $J,D,\mathcal B,\mathcal H_M,\mathcal H_E$ and
$\mathcal U$ from (\ref{eq:physical_tensor_bounds}) and its context.

A subspace $S\subseteq\mathcal H_M$ is $\delta$-viable for
$G\subseteq\mathcal H_M\otimes\mathcal H_E$ when
$\|P_{S\otimes\mathcal H_E}g\|^2\ge(1-\delta)\|g\|^2$ for every $g\in G$.
The closed spectral conditions and the requirement $\delta<1$ below are
recorded in \cite{gap:alvv16_spectral_projector_boundary}.
The norm estimate for the Chebyshev filter requires no positivity of its
unsquared polynomial~\cite{gap:alvv16_chebyshev_filter_positivity}.

\begin{theorem}[A common family for Chebyshev viability improvement]
    \label{thm:physical_chebyshev_viability}
    \lean{TNLean.MPS.AreaLaw.exists_common_physical_chebyshev_viable_family_fin}
    \leanok
    \uses{thm:physical_polynomial_tensor_family}
    There are a nonnegative integer $N\le D^2$ and matrices
    $W_0,\ldots,W_{N-1}$ on
    $\mathcal C(J)$, each with natural open-boundary bond bound
    $\mathcal B$, chosen before every supported exterior pair, every spectral
    choice and every input subspace below.
    After supplying $P_L,P_R$ supported on $\lbrack 0,a+1)$ and $\lbrack b+2\ell,L)$,
    define the actual Hamiltonian $H$ by
    (\ref{eq:physical_tensor_hamiltonian}) and put
    $H_m=\mathcal U H\mathcal U^*$.
    Suppose $H_m$ is self-adjoint and put
    $d=\dim(\mathcal H_M\otimes\mathcal H_E)$.
    Let $v_0,\ldots,v_{d-1}$ be its orthonormal eigenbasis, with eigenvalues
    $\lambda_j$.
    Supply real numbers $\eta_0<\eta_1\le u$. For every $0\le j<d$,
    suppose
    \begin{align}
        \lambda_j & \le u,\notag                              \\
        \lambda_j & \le\eta_0\ \text{or}\ \eta_1\le\lambda_j.
        \label{eq:physical_chebyshev_spectrum}
    \end{align}
    Define
    \begin{align}
        G      & =\spn_{\C}\{v_j:\lambda_j\le\eta_0\},
        \label{eq:physical_chebyshev_target}           \\
        \Delta & =\left(2\exp\left[-2k\sqrt{
            \frac{\eta_1-\eta_0}{u-\eta_0}}\right]\right)^2.
        \label{eq:physical_chebyshev_error}
    \end{align}
    For every subspace $S\subseteq\mathcal H_M$ and $\delta\in\R$ such
    that $S$ is $\delta$-viable for $G$ and $\delta<1$, put
    $V=\sum_{r=0}^{N-1}W_r(S)$. Then
    \begin{align}
        \dim V & \le D^2\dim S,\notag                                        \\
        V      & \text{ is }\frac{\Delta}{(1-\delta)^2}\text{-viable for }G.
        \label{eq:physical_chebyshev_conclusion}
    \end{align}
    These are supporting consequences for the spectral filtering and
    error-reduction arguments in \cite{Arad2016RigorousRG}.
\end{theorem}
\begin{proof}\leanok
    \uses{thm:physical_polynomial_tensor_family}
    Choose the common family in
    Theorem~\ref{thm:physical_polynomial_tensor_family}.
    The Chebyshev spectral estimate supplies a real polynomial $p$ of degree
    $k$ whose operator $K=p(H_m)$ preserves $G$ and $G^\perp$ and satisfies
    \begin{align}
        \|g\|    & \le\|Kg\|\quad(g\in G),\notag        \\
        \|Kw\|^2 & \le\Delta\|w\|^2\quad(w\in G^\perp).
        \label{eq:physical_chebyshev_contract}
    \end{align}
    Polynomial evaluation commutes with the matrix-to-operator identification
    and unitary conjugation. Applying
    (\ref{eq:physical_tensor_conjugation}) to the complex polynomial obtained
    from $p$ therefore gives $K=\sum_r W_r\otimes M_r$ with the same middle
    family. In particular, $K(S\otimes\mathcal H_E)\subseteq
        V\otimes\mathcal H_E$ and $\dim V\le N\dim S$.
    The two-factor error-reduction estimate applied to
    (\ref{eq:physical_chebyshev_contract}) gives viability
    $\Delta/(1-\delta)^2$ for $V$. Finally $N\le D^2$ gives the dimension
    bound in (\ref{eq:physical_chebyshev_conclusion}).
\end{proof}
